\documentclass{article}
\usepackage[utf8]{inputenc}
\usepackage{a4wide}

\begin{document}

\section*{เก็บน้ำฝน}

ดร.ต๋อย ต้องการเก็บน้ำฝนไว้ใช้งานยามหน้าแล้ง เลยจ้างคนงานมาขุดบ่อ แต่ด้วย ดร.ต๋อย

ต้องไปต่างประเทศนาน

ด้วยอะไรไม่ทราบทำให้คนงานจึงขุดบ่อเป็นหลุมเป็นบ่อขึ้นลงตามตัวอย่างรูปต่อไปนี้

กำหนดชุดจำนวนเต็มแทนความสูงของหลุมบ่อต่างๆ โดยหากความกว้างของแต่ละแท่งเท่ากับ 1

ถามว่า{\textbf{บ่อที่ขุดมานี้จะเก็บน้ำฝนได้เท่าใด}}\\



{\textbf{\hfill\break

}}



{\textbf{ตัวอย่างที่ 1}}



\begin{quote}

\textbf{Input:}



\textbf{10}



\textbf{1 0 2 1 0 1 2 1 2 1}\\



\textbf{Output: 6\\

}



{\textbf{\includegraphics[width=5.01042in,height=1.82292in,alt={Screenshot 2567-03-20 at 00.37.08.png}]{/public/upload/6e91b95280.png}\\

}}

\end{quote}



{\textbf{{\hfill\break

}}}



{\textbf{{ตัวอย่างที่ 2}\\

}}



\begin{quote}

\textbf{Input:}



\textbf{10}



\textbf{7 0 4 2 5 0 6 4 0 5}\\



\textbf{Output:25\\

}



\textbf{\includegraphics[width=3.30208in,height=1.82292in,alt={Screenshot 2567-03-20 at 00.38.26.png}]{/public/upload/da37707c10.png}\\

}

\end{quote}



\subsection*{Input}
บรรทัดแรกเป็นจำนวนของบ่อ



บรรทัดที่เหลือเป็นชุดเลขจำนวนเต็มแทนหมายเลขความสูงส่วนต่างๆ ในรูป (สีเข้ม) จำนวนไม่เกิน

300 ตัวแต่ละตัวมีค่า ระหว่าง 0 ถึง 100\\



\subsection*{Output}
ปริมาณน้ำฝนที่เก็บได้\\



\end{document}
